% Arithmetic contribution fragment. Reviewed source snapshot:
% openai/math adc7f1241b42e322a6451854ab7e4b4c146bf78a (accessed 2026-10-07).
% Claims attributed to the release are distinguished from established antecedents.

\subsection{Ordinary multiplicative correlations: retaining the terminal scale}
\label{sec:frontiers-ordinary-correlations}

Family 007 addresses the difference between cancellation at a specified cutoff
and cancellation after averaging over scales. Write $\lambda(n)=(-1)^{\Omega(n)}$
for the Liouville function, where $\Omega(n)$ counts prime factors with
multiplicity. Its manuscript states that an
absolute $c>0$ gives, for each fixed pair of nonproportional affine forms
$a_1n+b_1,a_2n+b_2$ with integer $a_i\geq1$ and $b_i\geq0$,
\[
 \left|\sum_{n\leq X}\lambda(a_1n+b_1)\lambda(a_2n+b_2)\right|
 \leq C_{a_1,a_2,b_1,b_2}\,X(\log X)^{-c},\qquad X\geq3.
\]
The coefficient-dependent constant can be ineffective; the statement fixes
the forms before $X$ grows \cite{arithOpenAI007}. The companion corrected
Elliott statement concerns unit-disc-valued multiplicative functions.
The obstruction is imitation of a periodic character multiplied by the
oscillatory factor $n^{it}$; the distance below measures this imitation
on primes. One fixed factor must satisfy, for every fixed Dirichlet character $\chi$,
$\inf_{|t|\leq N}\mathbb D(f,\chi n^{it};N)\to\infty$, where
$\mathbb D^2(f,g;N)=\sum_{p\leq N}(1-\Re(f(p)\overline{g(p)}))/p$.
Under this hypothesis the ordinary two-point mean tends to zero, including
fixed nonproportional affine forms. This second assertion supplies no rate
uniform over that function class \cite{arithOpenAI007}.

These are substantial changes of quantifiers relative to several established
results. Matom\"aki and Radziwi\l\l\ control multiplicative means in most
short intervals \cite{arithMR2016}; Matom\"aki, Radziwi\l\l\ and Tao prove
Chowla cancellation averaged over shifts \cite{arithMRT2015}. Tao's
two-point theorem controls logarithmically weighted affine correlations
\cite{arithTao2016}. Tao and Ter\"av\"ainen obtain ordinary correlations
outside an exceptional set of scales of logarithmic density zero
\cite{arithTT2019}. Helfgott--Radziwi\l\l\ divisibility graphs and
Pilatte's subsequent prime-product construction deliver increasingly strong
logarithmic cancellation \cite{arithHR2021,arithPilatte2026}. Pilatte already
uses composite steps, centred prime factors and nonbacktracking walks.
Consequently, the distinctive object of review is how the release combines
these ingredients to reach every cutoff with controlled error.

Two closer comparisons sharpen the remaining barrier. Klurman, Mangerel
and Ter\"av\"ainen obtain ordinary binary cancellation on a set of full
\emph{upper} logarithmic density when one fixed factor is nonpretentious
against every fixed character and fixed real twist
\cite[Definition~1.1 and Theorem~1.2]{arithKMT2023}.
Tao and Ter\"av\"ainen also obtain a logarithmic-power saving for slowly
growing affine Liouville forms outside a quantitatively controlled
exceptional set of scales \cite[Theorem~3.1 and Remarks~3.2]{arithTT2026}.
The release's quantitative assertion removes the exceptional cutoffs for
fixed forms; the latter comparison retains greater coefficient uniformity.
The release's qualitative assertion combines every-cutoff cancellation with the
stronger growing-twist hypothesis stated above. These are distinct changes
of scope and hypothesis. The prospective breakthrough is completing the
cutoff argument with its error budget intact, rather than obtaining an
ordinary average or logarithmic saving in isolation.

One unusual bridge is a finite-law comparison through circuit complexity.
The graph step uses products from disjoint prime supplies, with factors
$\mathbf1_{p\mid n}-1/p$. When the number of supplies grows like
$\log\log X$, the joint residue pattern repeats only after the product
of all participating primes, by the Chinese remainder theorem. That
period is too large for direct uniform sampling. The manuscript instead encodes residue observables by binary
variables and tests only bounded-depth circuits. Its finite-law lemma first
corrects small low-order biases to obtain an \emph{exactly} bounded-wise
independent distribution, then applies Braverman's circuit-fooling theorem.
The correction has the established almost-independence precedent of Alon,
Goldreich and Mansour \cite{arithAGM2003,arithBraverman2010}. Crucially,
translated sites retain the same residue for each prime. The comparison is
applied to the signed trace expansion before absolute values; treating
translated sites as independently resampled would change the arithmetic
model. The manuscript's finite-law, nonbacktracking-moment and
mixed-difference lemmas give separately inspectable checkpoints
\cite{arithOpenAI007}.

\begin{example}[Why translated sites share a prime residue]
Let $U$ be uniform modulo a prime $p$, and set
$c_0=\mathbf1_{U=0}-1/p$ and $c_h=\mathbf1_{U+h=0}-1/p$.
The same residue determines both translated sites, giving
\[
 \mathbb E[c_0c_h]=
 \begin{cases}
 (p-1)/p^2,&p\mid h,\\
 -1/p^2,&p\nmid h.
 \end{cases}
\]
Indeed, the two divisibility events coincide when $p\mid h$ and are
disjoint otherwise. Independently resampling the residue at the second
site would give zero in both cases, despite retaining each marginal
law. This small calculation explains why the manuscript's finite-law
comparison keeps shared residues inside the signed trace: independence
at the wrong level would erase the correlations being estimated.
\end{example}

The decisive quantitative bottleneck is uniformity of this comparison and
the high-trace count as the number of prime supplies grows, including all
squarefree padding choices. Repeated prime labels, rank witnesses and forest
encodings must produce a saving with absolute constants rather than costs
that erase the eventual logarithmic gain. A second, qualitative mechanism
fixes its graph scale before taking the averaging cutoff to infinity.
Its centre factors $\mathbf1_{p\mid n}-\theta/p$ need not have zero uniform
mean. Vertex normalization and signed singleton integration provide the
replacement cancellation; first-visit tree bookkeeping retains return
factors until averaging. This explains the order of limits and why the
qualitative theorem need not assume a rate of nonpretentiousness
\cite{arithOpenAI007}.

The documented formal targets cover both the quantitative fixed-form
assertion and the qualitative growing-twist assertion.\footnote{The
comparator configuration selects three final declarations:
\texttt{liouville\_log\_saving}, \texttt{binary\_corrected\_elliott} and
\texttt{affine\_corrected\_elliott}. The inspected Lean wrappers bind them to
the solution module \texttt{OAI.NumberTheory.TwoPointCorrelations.FinalMain}.
Their assumptions match the stated fixed-form and nonpretentiousness scope.
These declarations record the corresponding quantitative and qualitative
conclusions.}

\subsection{The exponent of $\pi$: logarithmic jets and determinant collisions}
\label{sec:frontiers-pi-exponent}

The irrationality exponent measures how well a real number can be
approximated by rationals: larger exponents allow infinitely many
unusually small approximation errors. Family 017 asserts the endpoint
$\mu(\pi)=2$. Its precise upper-bound statement is that, for every $\nu>2$, some
$Q(\nu)$ satisfies
\[
 |\pi-p/q|\geq q^{-\nu}\quad(p\in\mathbb Z,\ q\in\mathbb N,\ q\geq Q(\nu)).
\]
The threshold is ineffective. The assertion concerns $\pi$, not $\pi^2$, and
does not imply a fixed lower bound $c/q^2$ or bounded partial quotients
\cite{arithOpenAI017}. The contrast with established quantitative work is
sharp: Zeilberger and Zudilin's variant of Salikhov's integral method proves
the upper bound $7.103205334137\ldots$ \cite{arithZZ2020}. The released
manuscript seeks the endpoint through a different geometric and arithmetic
organization rather than optimizing those integrals.

Its central intermediate assertion is a separated-weight interpolation
theorem. A jet is a finite collection of Taylor coefficients, so jet
interpolation means prescribing values and derivatives at several
points. The question is whether enough polynomials remain available
after imposing the arithmetic constraints needed later in the proof.
Polynomials use a full weighted simplex degree budget; after
substituting $Y=1+t$ and $X_i=c_{ji}+u_i+\log(1+t)$, they must prescribe
specified jets at several coordinatewise distinct centres. The weights are
selected successively, \emph{before} the centres are known. The required
degree threshold may subsequently depend on the centres. This asymmetry is
essential: rational approximations determine centres later, so an argument
choosing weights from those approximations would not supply the needed
quantifier order \cite{arithOpenAI017}.

Dimension counting alone does not establish this surjectivity. The
manuscript proceeds through weighted transverse multiplicity, a curve-contact
inequality, and positivity on an ordinary blow-up of finite-colength jet
ideals. It then uses sufficiently large ordinary ideal powers and Serre
vanishing to produce the prescribed jets. Roth-type zero estimates on
products of commutative algebraic groups provide a relevant antecedent in
Farhi \cite{arithFarhi2006}; Demailly relates curve positivity, blow-ups and
asymptotic jet generation \cite{arithDemailly1992}. The release's proposed
bridge combines these themes for logarithmic curves with separated weights.
The decisive geometric checkpoint is its curve inequality and the resulting
centre-independent weight selection. The passage to ordinary ideal powers,
rather than integral closures, is another exact point requiring review
\cite{arithOpenAI017}.

The arithmetic half constructs a nonzero rational-Gaussian interpolation
minor whose denominator-clearing estimate supplies a lower bound. Translating
its rows expresses entries through entire exponential-polynomial functions.
Taylor degrees that coincide within a fixed transverse-index class create
equal rows and annihilate that determinant term. Taylor-degree cancellation
already underlies Laurent's interpolation-determinant estimates
\cite{arithLaurent2000}. Here the manuscript groups collisions by transverse
indices and combines them with full-simplex counts: with $m$ logarithmic
coordinates, the growth exponents in the degree parameter are $m+1$ for
the polynomial-space dimension and $m$ for the index-set cardinality.
Many low indices yield collision savings; few yield powers of small
approximation errors. The stated collision and determinant-comparison
lemmas expose the analytic upper bound that must contradict the arithmetic
lower bound \cite{arithOpenAI017}.

A concrete mathematical consequence would be convergence of the Flint Hills
series $\sum_{n\geq1}(n^3\sin^2 n)^{-1}$. Alekseyev proves that convergence
requires $\mu(\pi)\leq5/2$ \cite{arithAlekseyev2011}; Meiburg proves a
sufficient strict inequality $\mu(\pi)<5/2$ using sparsity of good rational
approximations \cite{arithMeiburg2022}. Thus the released exponent theorem
would settle this convergence question without a borderline argument. The
selected formal target records the approximation inequality and
exponent identity.\footnote{The comparator selects
\texttt{OAI.PiExponent.main}. A separate, unselected declaration in
\texttt{Main.lean} records Flint Hills summability.}
